%% ==== Résumé et Abstract ====
\chapter*{Résumé}
\addcontentsline{toc}{chapter}{Résumé}
\markboth{RÉSUMÉ}{RÉSUMÉ}
Ce mémoire compare l'algorithme EM classique et l'EM variationnel pour les modèles de mélange, en distinguant trois niveaux d'inférence : l'EM classique (paramètres ponctuels), l'EM variationnel à paramètres ponctuels et le Bayes variationnel (VB, paramètres aléatoires). Chaque avantage et chaque inconvénient avancé dans la littérature est relié à une preuve : théorème, expérience répliquée, expérience unique, ou mention « non testé », synthétisés dans une grille « axe $\times$ preuve ». Les algorithmes (EM, MAP-EM, VB, pLSA, LDA) sont implémentés en PyTorch avec un critère d'arrêt commun, et comparés à une baseline qui sélectionne $K$ par BIC ou ICL, sur dix scénarios synthétiques, un jeu semi-synthétique, un cas principal sur les régimes climatiques du Cameroun, un cas secondaire sur le risque de marché du S\&P 500 et une comparaison pLSA contre LDA sur un corpus réel.

La thèse défendue est que le VB ne « bat » pas l'EM : il déplace la sélection de $K$ vers le choix des hyperparamètres, apporte une incertitude approchée, au prix d'un coût et d'une sensibilité aux priors. Les résultats chiffrés sont ceux des tableaux générés par le code et de la grille du tableau 2.17 ; ils ne sont pas répétés ici.

\textbf{Mots-clés : }algorithme EM, EM variationnel, Bayes variationnel, mélanges gaussiens, sélection de modèle, ELBO, MAP-EM.

\chapter*{Abstract}
\addcontentsline{toc}{chapter}{Abstract}
\markboth{ABSTRACT}{ABSTRACT}
This thesis compares the classical EM algorithm with variational EM for mixture models, distinguishing three levels of inference: classical EM (point parameters), variational EM with point parameters, and variational Bayes (VB, random parameters). Every advantage and drawback reported in the literature is tied to evidence: a theorem, a replicated experiment, a single experiment, or the label ``untested'', summarised in an ``axis $\times$ evidence'' grid. The algorithms (EM, MAP-EM, VB, pLSA, LDA) are implemented in PyTorch with a common stopping rule and compared with a baseline that selects $K$ by BIC or ICL, on ten synthetic scenarios, a semi-synthetic benchmark, a main case study on the climate regimes of Cameroon, a secondary case on S\&P 500 market risk, and a pLSA versus LDA comparison on a real corpus.

The thesis defended is that VB does not ``beat'' EM: it shifts the choice of $K$ towards the choice of hyperparameters and provides an approximate uncertainty, at the price of a higher cost and a sensitivity to priors. Numerical results are those of the tables generated by the code and of the grid in Table 2.17; they are not repeated here.

\textbf{Keywords: }EM algorithm, variational EM, variational Bayes, Gaussian mixtures, model selection, ELBO, MAP-EM.
